\documentclass[11pt,a4paper]{article}

\usepackage[utf8]{inputenc}
\usepackage[T1]{fontenc}
\usepackage[spanish]{babel}
\usepackage{amsmath,amssymb}
\usepackage{graphicx}
\usepackage{geometry}
\geometry{margin=2.5cm}

\title{Laboratorio N° 01 --- Métodos Matemáticos para Físicos I (CF2C2)\\
\large Borrador de informe}
\author{Universidad Nacional de Ingeniería --- Facultad de Ciencias\\
Escuela Profesional de Física --- Ciclo 2026-I}
\date{Fecha de entrega: viernes 17 de abril de 2026}

\begin{document}
\maketitle

\section*{Instrucciones del laboratorio (resumen)}
Este laboratorio es de carácter domiciliario. Cada problema debe resolverse tanto analíticamente (mostrando todos los pasos matemáticos) como computacionalmente usando Python (SymPy, NumPy, Matplotlib) o Wolfram Mathematica. El informe en PDF debe incluir el desarrollo analítico y las gráficas generadas por computadora. Los códigos deben ir en archivos separados y citarse en el informe. Justificar rigurosamente cada paso.

\vspace{2em}

% ---------------------------------------------------------------------------
\section*{Problema 1. Ecuaciones lineales de coeficientes constantes no homogéneas}

\subsection*{Enunciado}
\textbf{(a) (3p)} Resuelva analíticamente las siguientes EDOs usando reducción de orden (sustitución $y' = v$), y verifique su solución con \texttt{dsolve} de SymPy o \texttt{DSolve} de Mathematica:
\begin{enumerate}
  \item[(i)] $y'' + 7y' = e^{-7x}$,
  \item[(ii)] $y'' - 7y' = (x - 1)^2$.
\end{enumerate}

\textbf{(b) (2p)} Usando Python o Mathematica, grafique en una misma figura las familias de soluciones de ambas EDOs para al menos cuatro pares de condiciones iniciales $(y(0), y'(0))$ de su elección. Comente brevemente el comportamiento cualitativo de cada familia.

\noindent\textit{Cumplimiento:} el apartado (a) exige la reducción de orden en el desarrollo a mano y la verificación simbólica (\texttt{dsolve}); el (b) exige una sola figura con ambas EDOs, al menos cuatro pares de CI y comentario cualitativo.

\subsection*{Desarrollo analítico}
En la página siguiente se muestra la resolución analítica a mano del problema~1.
\clearpage
\begin{figure}[p]
  \centering
  \includegraphics[width=\linewidth,height=0.92\textheight,keepaspectratio]{images/analitico/p1.jpg}
  \caption{Problema~1: desarrollo analítico a mano.}
  \label{fig:analitico-p1}
\end{figure}
\clearpage

\subsection*{Verificación con SymPy (\texttt{dsolve}), apartado (a)}
Ejecutando \texttt{1\_de\_simbolic.py}, \texttt{sympy.dsolve} devuelve las soluciones generales siguientes (constantes de integración $C_1$, $C_2$ en la convención de SymPy). Deben coincidir---salvo posibles reparametrizaciones de constantes---con el resultado obtenido por reducción de orden en la figura~\ref{fig:analitico-p1}.

\textbf{(i) $y''+7y'=e^{-7x}$.}
\begin{equation}
  y(x)=C_1+\left(C_2-\frac{x}{7}\right)\mathrm{e}^{-7x}
  \;=\; C_1+C_2\,\mathrm{e}^{-7x}-\frac{x}{7}\,\mathrm{e}^{-7x}.
\end{equation}

\textbf{(ii) $y''-7y'=(x-1)^2$.}
\begin{equation}
  y(x)=C_1+C_2\,\mathrm{e}^{7x}-\frac{x^{3}}{21}+\frac{6x^{2}}{49}-\frac{37x}{343}.
\end{equation}

\subsection*{Parte computacional, gráficas y comentario}
La figura~\ref{fig:lab1} reúne en un solo lienzo los dos paneles pedidos (familias con cuatro pares $(y(0),y'(0))$). Cualitativamente, en el panel izquierdo la exponencial $e^{-7x}$ en el término fuente modula la respuesta frente a la homogénea con raíz cero; en el derecho, al cambiar el signo de $y'$, cambia el factor creciente/decayente de la homogénea ($e^{7x}$ frente a constantes en el otro caso).

\begin{figure}[htbp]
  \centering
  \includegraphics[width=\linewidth]{images/lab1_familias_edos.pdf}
  \caption{Problema~1(b): familias de soluciones para ambas EDOs (Python/SymPy/Matplotlib).}
  \label{fig:lab1}
\end{figure}

\subsection*{Referencias a archivos de código}
El código va en archivo separado y se cita aquí según indica el laboratorio:
\begin{itemize}
  \item \texttt{1\_de\_simbolic.py}: imprime con \texttt{sympy.dsolve} las soluciones generales de (i)--(ii) para contrastarlas con el cálculo a mano, y genera la figura~\ref{fig:lab1} imponiendo $(y(0),y'(0))$ en $x=0$.
\end{itemize}

\newpage

% ---------------------------------------------------------------------------
\section*{Problema 2. Ecuación de Euler y sustitución $x = e^t$}

\subsection*{Enunciado}
Resuelva analíticamente la siguiente ecuación diferencial mediante la sustitución $x = e^t$:
\[
  x^2 y'' + x y' - y = x^m, \qquad |m| \neq 1,\quad x > 0.
\]

\textbf{(a) (2p)} Transforme la ecuación a una EDO lineal de coeficientes constantes en la variable $t$ y resuélvala completamente (considere los tres casos posibles para $m$: $m \neq \pm 1$, con $m$ real).

\textbf{(b) (1p)} Regrese a la variable original $x$ y escriba la solución general $y(x)$.

\textbf{(c) (1p)} Con Python o Mathematica, grafique $y(x)$ en $[0{,}1,\, 4]$ para los casos particulares $m = 0$, $m = 2$ y $m = 3$ (con constantes de integración $c_1 = c_2 = 1$). Muestre las tres curvas en una misma figura.

\noindent\textit{Cumplimiento:} los apartados (a)--(b) son analíticos (sustitución $x=e^t$, casos de $m$); el (c) es la gráfica única con $m\in\{0,2,3\}$ en $[0{,}1,4]$.

\subsection*{Desarrollo analítico}
En la página siguiente se muestra la resolución analítica a mano del problema~2.
\clearpage
\begin{figure}[p]
  \centering
  \includegraphics[width=\linewidth,height=0.92\textheight,keepaspectratio]{images/analitico/p2.jpg}
  \caption{Problema~2: desarrollo analítico a mano.}
  \label{fig:analitico-p2}
\end{figure}
\clearpage

\subsection*{Parte computacional y gráficas}
\begin{figure}[htbp]
  \centering
  \includegraphics[width=\linewidth]{images/lab2_euler_c.pdf}
  \caption{Problema~2(c): $y(x)$ para $m=0,2,3$ con $c_1=c_2=1$ en $[0{,}1,4]$ (una figura, tres curvas).}
  \label{fig:lab2}
\end{figure}

\subsection*{Referencias a archivos de código}
\begin{itemize}
  \item \texttt{2\_euler.py}: evalúa la solución general adoptada y genera la figura~\ref{fig:lab2}.
\end{itemize}

\newpage

% ---------------------------------------------------------------------------
\section*{Problema 3. Ecuación de Euler no homogénea}

\subsection*{Enunciado}
\textbf{(a) (3p)} Resuelva analíticamente el siguiente problema de valor inicial:
\[
  x^2 y'' - 2x y' + 2y = x^2 - 2x + 2, \qquad y(1) = 1,\quad y'(1) = 0.
\]

\textbf{(b) (3p)} Usando Python o Mathematica, resuelva numéricamente el mismo problema en el intervalo $[0{,}5,\, 5]$ y compare la solución numérica con la analítica en una gráfica. Calcule y reporte el error máximo absoluto $\max_x |y_{\mathrm{num}}(x) - y_{\mathrm{anal}}(x)|$.

\noindent\textit{Cumplimiento:} (a) es el PVI a mano; (b) exige integración numérica en $[0{,}5,5]$, gráfica comparativa y valor de $\max_x|y_{\mathrm{num}}-y_{\mathrm{anal}}|$ (en la práctica sobre una malla fina; ver script citado).

\subsection*{Desarrollo analítico}
En la página siguiente se muestra la resolución analítica a mano del problema~3.
\clearpage
\begin{figure}[p]
  \centering
  \includegraphics[width=\linewidth,height=0.92\textheight,keepaspectratio]{images/analitico/p3.jpg}
  \caption{Problema~3: desarrollo analítico a mano.}
  \label{fig:analitico-p3}
\end{figure}
\clearpage

\subsection*{Parte computacional, comparación y error máximo}
\begin{figure}[htbp]
  \centering
  \includegraphics[width=\linewidth]{images/lab3_numerico_vs_analitico.pdf}
  \caption{Problema~3(b): solución numérica (SciPy) frente a analítica (SymPy) en $[0{,}5,5]$.}
  \label{fig:lab3}
\end{figure}

Al ejecutar \texttt{3\_numerical.py} se imprime el error máximo absoluto sobre la malla usada en el script (RK45 con \texttt{dense\_output}, tolerancias del propio archivo). Ejemplo de orden de magnitud obtenido en una corrida típica: $E_{\max}\sim 10^{-9}$.

\subsection*{Referencias a archivos de código}
\begin{itemize}
  \item \texttt{3\_numerical.py}: \texttt{sympy.dsolve} con las CI del enunciado, \texttt{scipy.integrate.solve\_ivp} desde $x=1$ hacia los extremos del intervalo, figura~\ref{fig:lab3} y cálculo de $E_{\max}$.
\end{itemize}

\newpage

% ---------------------------------------------------------------------------
\section*{Problema 4. Operadores diferenciales lineales y el Wronskiano}

\subsection*{Enunciado}
\textbf{(a) (2p)} Dado el operador diferencial lineal $L = x^3 D^3 + (\sin x)\, D^2 + e^x D$, donde $D = \mathrm{d}/\mathrm{d}x$, determine el dominio del operador. Luego, aplique $L$ a las funciones $f(x) = \cos x$ y $g(x) = x^2$ y exprese explícitamente $(Lf)(x)$ y $(Lg)(x)$.

\textbf{(b)} Considere las tres funciones: $y_1 = e^x$, $y_2 = x e^x$, $y_3 = x^2 e^x$.
\begin{enumerate}
  \item[\textbf{i.}] \textbf{(1p)} Calcule analíticamente el Wronskiano $W[y_1, y_2, y_3](x)$.
  \item[\textbf{ii.}] \textbf{(1p)} Determine una EDO lineal homogénea de tercer orden con coeficientes constantes de la cual $\{y_1, y_2, y_3\}$ sea el conjunto fundamental de soluciones.
  \item[\textbf{iii.}] \textbf{(1p)} Verifique su Wronskiano con Python (usando \texttt{sympy.wronskian}) y gráfiquelo en $[-2, 2]$.
\end{enumerate}

\noindent\textit{Cumplimiento:} el apartado (a) y el (b)\,i--ii son enteramente analíticos (dominio de $L$, aplicación a $f,g$, Wronskiano y EDO homogénea); solo el (b)\,iii exige código Python con \texttt{sympy.wronskian} y gráfica en $[-2,2]$.

\subsection*{Desarrollo analítico}
En la página siguiente se muestra la resolución analítica a mano del problema~4 (primera hoja).
\clearpage
\begin{figure}[p]
  \centering
  \includegraphics[width=\linewidth,height=0.92\textheight,keepaspectratio]{images/analitico/p4_1.jpg}
  \caption{Problema~4: desarrollo analítico a mano (hoja~1).}
  \label{fig:analitico-p4a}
\end{figure}
\clearpage

En la página siguiente se muestra la resolución analítica a mano del problema~4 (segunda hoja).
\clearpage
\begin{figure}[p]
  \centering
  \includegraphics[width=\linewidth,height=0.92\textheight,keepaspectratio]{images/analitico/p4_2.jpg}
  \caption{Problema~4: desarrollo analítico a mano (hoja~2).}
  \label{fig:analitico-p4b}
\end{figure}
\clearpage

\subsection*{Parte computacional y gráficas}
En el desarrollo analítico (figuras~\ref{fig:analitico-p4a}--\ref{fig:analitico-p4b}) se obtiene $W[y_1,y_2,y_3](x)=2\mathrm{e}^{3x}$. El script compara esa expresión con la que devuelve \texttt{sympy.wronskian}: la diferencia simbólica es nula y, en la malla de la gráfica, el máximo de $|W_{\mathrm{SymPy}}-W_{\mathrm{manuscrito}}|$ resulta $0$ (salida de consola al ejecutar \texttt{4\_wronskian.py}).

\begin{figure}[htbp]
  \centering
  \includegraphics[width=\linewidth]{images/lab4_wronskiano.pdf}
  \caption{Problema~4(b)\,iii: superposición de $W$ vía \texttt{sympy.wronskian} y de $2\mathrm{e}^{3x}$ (analítico a mano) en $[-2,2]$; las curvas coinciden.}
  \label{fig:lab4}
\end{figure}

\subsection*{Referencias a archivos de código}
\begin{itemize}
  \item \texttt{4\_wronskian.py}: calcula $W$ con \texttt{sympy.wronskian}, lo confronta con $2\mathrm{e}^{3x}$, imprime la diferencia simbólica y el error máximo en malla, y genera la figura~\ref{fig:lab4}.
\end{itemize}

\end{document}
